% Part: applied-modal-logics
% Chapter: epistemic-logic
% Section: relational-models

\documentclass[../../../include/open-logic-section]{subfiles}

\begin{document}

\olfileid{aml}{el}{rel}

\olsection{النماذج العلائقية}

أصل الدلالة في منطق المعرفة هو أصلها في المنطق الجهوي المعتاد: نموذج
علائقي فيه عوالم، وإسناد يعين عند أي عالم تصدق ال!!{propositional variable}s.
فإن كان الفاعل واحدًا كانت علاقة النفاذ واحدة؛ وإن تعدد الفاعلون تعددت
علاقات النفاذ، فلكل~$a$ من مجموعة الرموز~$G$ علاقة.

فــ\emph{النموذج العلائقي}، إذن، مجموعة عوالم تصل بينها علاقات نفاذ
ثنائية بعدد الفاعلين، ومعها إسناد يعين ال!!{propositional variable}s
الصادقة عند كل عالم.

\begin{defn}
  \emph{نموذج} لغة المعرفة المتعددة الفاعلين هو الثلاثي
  $\mModel{M}=\tuple{W,(R_a)_{a \in G},V}$، وفيه:
  \begin{enumerate}
  \item $W$ مجموعة غير فارغة من «العوالم»؛
  \item لكل~$a \in G$، تكون~${R}_a$ علاقة نفاذ ثنائية على~$W$؛
  \item $V$ دالة تسند إلى كل~!!{propositional
    variable}~$p$ مجموعة~$V(p)$ من العوالم الممكنة.
  \end{enumerate}
  فإذا كان~$R_a ww'$ قلنا إن الفاعل~$a$ \emph{يمكنه النفاذ إلى}~$w'$
  من~$w$؛ وإذا كان~$w \in V(p)$ قلنا إن~$p$ \emph{صادق عند}~$w$.
\end{defn}

فليس وراء ذلك إلا آلية المنطق الجهوي النظامي مع تكثير علاقات النفاذ.
ونحمل علاقة كل فاعل على حاله المعلوماتية: فقولنا~$R_a ww'$ يعني، في
التفسير المعتاد، أن~$w'$ متسق مع معلومات~$a$ عند~$w$؛ أي إن~$a$ عند
$w$ لا يميز العالم~$w$ من العالم~$w'$.

\end{document}
